%% -------------------------------------------------------
%% Kapitel 7: Testing und Deployment
%% -------------------------------------------------------
\chapter{Testing und Deployment}
\label{chap:testing}

\section{Teststrategie}
\label{sec:teststrategie}

Eine systematische Teststrategie stellt sicher, dass das Backend und die iOS-App korrekt
funktionieren und die definierten Anforderungen erfüllen. Für DoggoGO wurden drei
Testebenen eingesetzt:

\begin{enumerate}
  \item \textbf{Manuelle API-Tests} über Swagger~UI und Postman
  \item \textbf{Unit-Tests} für einzelne Backend-Komponenten (xUnit)
  \item \textbf{Integrationstests} für das Zusammenspiel zwischen iOS-App und Backend
\end{enumerate}

\section{Backend-Tests}
\label{sec:backend-tests}

\subsection{Manuelle Tests mit Swagger und Postman}
\label{subsec:swagger-tests}

Während der Entwicklung wurden alle API-Endpunkte zunächst manuell über die
\textbf{Swagger-UI} getestet. Dies erlaubt schnelles Feedback ohne Testrahmenaufwand:
Registrierung, Anmeldung, CRUD-Operationen für Hunde sowie das Abrufen der Rassenliste
wurden iterativ überprüft. Für komplexere Szenarien (z.\,B.\ Fehlerfälle, ungültige Tokens)
wurde \textbf{Postman} mit einer gespeicherten Collection genutzt, die die wichtigsten
Endpunkte mit vorbereiteten Request-Bodies enthält.

Tabelle~\ref{tab:test-szenarien} listet die wichtigsten getesteten Szenarien auf:

\begin{table}[H]
  \centering
  \caption{Getestete API-Szenarien}
  \label{tab:test-szenarien}
  \small
  \begin{tabularx}{\textwidth}{|l|l|Y|l|}
    \hline
    \textbf{Szenario} & \textbf{Endpunkt} & \textbf{Erwartetes Ergebnis} &
    \textbf{Status} \\
    \hline
    Registrierung (gültig)     & POST /owners/register & 201 Created, Owner-Objekt    & OK \\
    \hline
    Registrierung (doppelt)    & POST /owners/register & 409 Conflict (doppelte Email) & OK \\
    \hline
    Anmeldung (gültig)         & POST /owners/login    & 200 OK, JWT                  & OK \\
    \hline
    Anmeldung (Falsches PW)    & POST /owners/login    & 401 Unauthorized             & OK \\
    \hline
    Hund anlegen (auth.)       & POST /api/dogs        & 201 Created, Dog-Objekt      & OK \\
    \hline
    Hund anlegen (unauth.)     & POST /api/dogs        & 401 Unauthorized             & OK \\
    \hline
    Hund löschen (Cascade)     & DELETE /owners/\{id\} & 204 + Hunde gelöscht         & OK \\
    \hline
    Rassen abrufen             & GET /api/breeds       & 200 OK, Rassenliste mit Regeln & OK \\
    \hline
  \end{tabularx}
\end{table}

\subsection{Unit-Tests mit xUnit}
\label{subsec:xunit}

Für die Passwort-Hashing-Logik und die JWT-Erzeugung wurden Unit-Tests mit
\textbf{xUnit} implementiert. Listing~\ref{lst:unittest} zeigt einen exemplarischen Test
für den Passwort-Hash-Vergleich:

\begin{lstlisting}[language=csharp, caption={xUnit-Test – Passwort-Hash-Verifikation}, label=lst:unittest, float=htbp]
public class PasswordHashTests
{
    [Fact]
    public void VerifyPassword_CorrectPassword_ReturnsTrue()
    {
        // Arrange
        using var hmac = new HMACSHA512();
        var salt = hmac.Key;
        var hash = hmac.ComputeHash(Encoding.UTF8.GetBytes("Test1234!"));

        // Act
        using var hmac2 = new HMACSHA512(salt);
        var verifyHash = hmac2.ComputeHash(
            Encoding.UTF8.GetBytes("Test1234!"));

        // Assert
        Assert.Equal(hash, verifyHash);
    }

    [Fact]
    public void VerifyPassword_WrongPassword_ReturnsFalse()
    {
        using var hmac = new HMACSHA512();
        var salt = hmac.Key;
        var hash = hmac.ComputeHash(Encoding.UTF8.GetBytes("Test1234!"));

        using var hmac2 = new HMACSHA512(salt);
        var verifyHash = hmac2.ComputeHash(
            Encoding.UTF8.GetBytes("WrongPassword"));

        Assert.NotEqual(hash, verifyHash);
    }
}
\end{lstlisting}

\section{iOS-Tests}
\label{sec:ios-tests}

Die iOS-App wurde mit \textbf{XCTest} getestet. Aufgrund des Umfangs der Diplomarbeit
konzentrierten sich die Tests auf das \texttt{AuthService}-Modul (Keychain-Operationen)
und die \texttt{APIService}-Klasse (URL-Konstruktion, Fehlerbehandlung):

\begin{lstlisting}[language=swift, caption={XCTest – AuthService Keychain-Test}, label=lst:xctest, float=htbp]
class AuthServiceTests: XCTestCase {

    func testSaveAndRetrieveToken() {
        let service = AuthService.shared
        let testToken = "eyJhbGciOiJIUzI1NiIsInR5cCI6IkpXVCJ9.test"

        service.saveToken(testToken)
        let retrieved = service.getToken()

        XCTAssertEqual(testToken, retrieved)
    }

    func testDeleteToken() {
        let service = AuthService.shared
        service.saveToken("some.token.value")
        service.deleteToken()

        XCTAssertNil(service.getToken())
    }
}
\end{lstlisting}

\section{Integrationstests}
\label{sec:integrationstests}

Integrationstests überprüften das Ende-zu-Ende-Verhalten des Gesamtsystems: Die iOS-App
wird im Simulator gestartet, und die Benutzerabläufe (Registrierung → Anmeldung →
Hundeprofi anlegen → Rasse mit Rechtsvorschriften abrufen) werden manuell durchgeführt.
Die HTTP-Anfragen wurden dabei mit dem Netzwerk-Inspector in Xcode beobachtet, um
korrekte Request/Response-Zyklen zu verifizieren.

\section{Deployment-Verifikation}
\label{sec:deployment-test}

Nach einem \texttt{docker compose up --build} wurden folgende Punkte geprüft:

\begin{itemize}
  \item PostgreSQL-Container läuft und ist über Port~5432 erreichbar.
  \item Backend-Container startet, Datenbankmigrationen werden ohne Fehler angewendet.
  \item Breed-Seeding protokolliert \textit{„Breeds seeded successfully"}.
  \item Swagger-UI ist unter \texttt{http://localhost:8080/swagger} erreichbar.
  \item Registrierung und Anmeldung über Swagger-UI funktionieren korrekt.
\end{itemize}
